\documentclass[12pt]{article}

\usepackage[margin=1.12in]{geometry}
\usepackage{amsmath,amssymb,amsthm,mathtools}
\usepackage{booktabs}
\usepackage{hyperref}

\newtheorem{theorem}{Theorem}
\newtheorem{proposition}{Proposition}
\theoremstyle{remark}
\newtheorem{remark}{Remark}

\setlength{\emergencystretch}{3em}
\hypersetup{
  pdftitle={Discrete Black-Hole Spectroscopy Does Not Force Binary Rank Without an Alphabet Law},
  pdfsubject={Bekenstein-Mukhanov spectroscopy, finite alphabets, black-hole entropy},
  pdfkeywords={black-hole spectroscopy, Bekenstein-Mukhanov, finite alphabet, binary rank, area quantization}
}

\title{Discrete Black-Hole Spectroscopy Does Not Force\\
Binary Rank Without an Alphabet Law}
\author{}
\date{Live review note of June 20, 2026}

\begin{document}
\maketitle

\begin{abstract}
Bekenstein--Mukhanov spectroscopy is a genuine possible positive route to a
finite alphabet: under an integer degeneracy law \(g_n=k^n\), a resolved
fundamental area spacing \(\Delta A=4\log(k)\,\ell_P^2\) in fixed units would
determine \(k\), and \(\Delta A=4\log(2)\,\ell_P^2\) would determine \(k=2\).
The point of this note is the converse boundary.  Entropy, qualitative
discreteness, recovery data, or an area spacing without the physical law
identifying that spacing with an integer degeneracy base do not force the
binary predicate \(q=2\).  Binary rank descends from spectroscopy exactly when
the supplied spectral law and measured data are injective on the
binary-versus-nonbinary comparison fiber.
\end{abstract}

\paragraph{Status.}
This is non-frozen external-review metadata for Team 1.  It is intended as a
small claim that a black-hole spectroscopy or quantum-gravity reviewer can
classify as known, false, missing-hypothesis, or useful.  It does not modify
any frozen Team 1 manuscript, Lean archive, public mirror manifest, or SHA-256
package.

\section{Primary sources}

The boundary below was checked against Bekenstein's entropy-area proposal
\cite{BekensteinEntropy}, Bekenstein--Mukhanov quantum black-hole spectroscopy
\cite{BekensteinMukhanov}, Bekenstein's black-hole-as-atom review
\cite{BekensteinAtoms}, and Maggiore's quasinormal-mode interpretation of
area quantization \cite{MaggioreQNM}.  These works are used as positive
spectroscopy or area-quantization inputs, not as claims that binary rank is
absent from nature.

\section{Alphabet identifiability}

Let \(q\ge2\) be the proposed finite alphabet size of a horizon cell or finite
central primitive.  Let \(D(q)\) denote the actual spectroscopy data reported
under a specified physical law: resolved fundamental frequency or area
spacing, transition intensities, level degeneracies, Planck/area units, and
the rule connecting those data to \(q\).

\begin{theorem}[spectroscopy alphabet criterion]
\label{thm:spectral-alphabet}
The binary predicate \(q=2\) is reconstructible from the supplied spectroscopy
data \(D\) if and only if, inside the declared comparison class,
\[
  D(q)=D(2)\quad\Longrightarrow\quad q=2 .
  \tag{1}
\]
Equivalently, the spectroscopy map must be injective on the
binary-versus-nonbinary fiber.
\end{theorem}

\begin{proof}
If \(q=2\) is reconstructible from \(D\), then the predicate is a function of
the data, so any model with the same data as \(q=2\) must also satisfy
\(q=2\).  Conversely, if (1) holds, define the binary predicate on the image of
\(D\) by declaring the fiber \(D(2)\) binary and every other image nonbinary.
Condition (1) is exactly the well-definedness condition for this assignment on
the relevant comparison class.
\end{proof}

\section{Bekenstein--Mukhanov as a positive exit}

In the Bekenstein--Mukhanov idealization, the area levels are equally spaced
and the degeneracy of the \(n\)-th level has the form
\[
  g_n=k^n,\qquad k\in\{2,3,\ldots\}.
\]
Matching the entropy \(S_n=\log g_n=n\log k\) to the
Bekenstein--Hawking normalization \(S=A/(4\ell_P^2)\) gives
\[
  A_n=4n\log(k)\,\ell_P^2,
  \qquad
  \Delta A=4\log(k)\,\ell_P^2 .
  \tag{2}
\]

\begin{proposition}[what the line measurement supplies]
\label{prop:bm-positive}
Assume the Bekenstein--Mukhanov integer-degeneracy law \(g_n=k^n\), fixed
Planck/area units, and a resolved fundamental spacing \(\Delta A\).  Then the
spacing determines \(k\) exactly when
\[
  k=\exp\!\left(\frac{\Delta A}{4\ell_P^2}\right)
\]
is a positive integer satisfying the declared law.  In particular,
\(\Delta A=4\log(2)\,\ell_P^2\) gives \(k=2\).
\end{proposition}

\begin{proof}
Equation (2) gives \(\Delta A=4\log(k)\ell_P^2\).  Solving for \(k\) gives the
displayed expression.  The conclusion that \(k\) is an alphabet size uses the
integer-degeneracy law; the measured spacing alone supplies only the numerical
spacing parameter.
\end{proof}

\section{What does not force binary rank}

\begin{proposition}[spacing without an alphabet law is nonidentifying]
\label{prop:nonidentifying}
The following data do not by themselves imply \(q=2\):
\begin{enumerate}
\item entropy proportional to area;
\item qualitative discreteness of the area or mass spectrum;
\item a uniformly spaced area spectrum \(A_n=\gamma n\ell_P^2\) with
      unspecified \(\gamma\);
\item a spacing parameter \(\gamma\) without a law identifying
      \(\gamma=4\log(q)\) for an integer finite alphabet \(q\);
\item recovery, sector, modular, crossed-product, or quotient-visible records
      that do not contain a \(q\)-separating spectral measurement.
\end{enumerate}
\end{proposition}

\begin{proof}
Entropy gives a number of states, \(\exp(S)\), not a distinguished base for
writing that number as powers of \(2\), \(3\), or another alphabet.  A
uniformly spaced spectrum with free parameter \(\gamma\) is compatible with
\(\gamma=4\log(2)\), \(\gamma=4\log(3)\), \(\gamma=8\pi\), or other proposed
area quanta, depending on the physical quantization rule.  Therefore the data
listed above fail condition (1): the binary and nonbinary comparison classes
are not separated until the finite-alphabet law or direct measurement channel
is part of the supplied data.
\end{proof}

\section{Referee question}

The exact external question is:

\begin{quote}
Does the black-hole spectroscopy or area-quantization hypothesis in the target
setting supply a resolved finite-alphabet law whose data are injective on
\(q=2\), or does it supply only entropy, qualitative discreteness, or area
spacing without binary-alphabet accountability?
\end{quote}

If the first holds, spectroscopy is a positive exit from the Team 1 finite
central obstruction.  If the second holds, Team 1 should not claim a
spectroscopic derivation of binary rank; the correct conclusion is only that a
direct finite-alphabet measurement would be enough.

\begin{thebibliography}{9}

\bibitem{BekensteinEntropy}
J. D. Bekenstein,
\emph{Black Holes and Entropy},
Phys. Rev. D 7 (1973), 2333,
\href{https://link.aps.org/doi/10.1103/PhysRevD.7.2333}{doi:10.1103/PhysRevD.7.2333}.

\bibitem{BekensteinMukhanov}
J. D. Bekenstein and V. F. Mukhanov,
\emph{Spectroscopy of the quantum black hole},
Phys. Lett. B 360 (1995), 7--12,
\href{https://arxiv.org/abs/gr-qc/9505012}{arXiv:gr-qc/9505012}.

\bibitem{BekensteinAtoms}
J. D. Bekenstein,
\emph{Quantum Black Holes as Atoms},
in Proceedings of the Eighth Marcel Grossmann Meeting,
\href{https://arxiv.org/abs/gr-qc/9710076}{arXiv:gr-qc/9710076}.

\bibitem{MaggioreQNM}
M. Maggiore,
\emph{The physical interpretation of the spectrum of black hole quasinormal modes},
Phys. Rev. Lett. 100 (2008), 141301,
\href{https://arxiv.org/abs/0711.3145}{arXiv:0711.3145}.

\end{thebibliography}

\end{document}
